\documentclass[12pt, a4paper]{article}
\usepackage[utf8]{inputenc}
\usepackage[ngerman]{babel}
\usepackage{amsmath, amssymb}
\usepackage{geometry}
\geometry{margin=2.5cm}
\usepackage{parskip}

\title{\textbf{Aufgabe 6 -- Lösung} \\ Methoden I: Mathematik, 2026S}
\author{}
\date{}

\begin{document}
\maketitle

\section*{Benötigte Formeln}

\begin{align}
    \sum_{k=1}^{n} 1 &= n \\[6pt]
    \sum_{k=1}^{n} k &= \frac{n(n+1)}{2} \\[6pt]
    \sum_{k=1}^{n} k^2 &= \frac{n(n+1)(2n+1)}{6} \\[6pt]
    \sum_{k=1}^{n} k^3 &= \left(\frac{n(n+1)}{2}\right)^2
\end{align}

\textbf{Wichtige Rechenregel:} Konstante Faktoren und Summen kann man aufteilen:
\[
    \sum_{k=1}^{n} (a \cdot f(k) + b \cdot g(k)) = a \cdot \sum_{k=1}^{n} f(k) + b \cdot \sum_{k=1}^{n} g(k)
\]

% ============================================================
\section*{(a) $\displaystyle\sum_{k=1}^{42} \left(k^2 - \frac{k}{2} - 13\right)$}

Aufteilen in drei separate Summen:

\[
    \sum_{k=1}^{42} k^2 - \frac{1}{2} \sum_{k=1}^{42} k - 13 \sum_{k=1}^{42} 1
\]

Jetzt einzeln berechnen mit den Formeln (jeweils $n = 42$):

\bigskip

\textbf{Summe 1:} $\displaystyle\sum_{k=1}^{42} k^2 = \frac{42 \cdot 43 \cdot 85}{6} = \frac{153510}{6} = 25585$

\textit{Erklärung: $2n + 1 = 2 \cdot 42 + 1 = 85$, dann $42 \cdot 43 = 1806$, dann $1806 \cdot 85 = 153510$, geteilt durch $6$.}

\bigskip

\textbf{Summe 2:} $\displaystyle\frac{1}{2} \sum_{k=1}^{42} k = \frac{1}{2} \cdot \frac{42 \cdot 43}{2} = \frac{1}{2} \cdot 903 = 451{,}5$

\bigskip

\textbf{Summe 3:} $\displaystyle 13 \sum_{k=1}^{42} 1 = 13 \cdot 42 = 546$

\bigskip

\textbf{Ergebnis:}
\[
    25585 - 451{,}5 - 546 = 24587{,}5
\]

\[
    \boxed{\sum_{k=1}^{42} \left(k^2 - \frac{k}{2} - 13\right) = 24587{,}5 = \frac{49175}{2}}
\]

% ============================================================
\section*{(b) $\displaystyle\sum_{k=1}^{21} \left(2k^3 - 2k + 17\right)$}

Aufteilen:

\[
    2 \sum_{k=1}^{21} k^3 - 2 \sum_{k=1}^{21} k + 17 \sum_{k=1}^{21} 1
\]

\bigskip

\textbf{Summe 1:} $\displaystyle 2 \sum_{k=1}^{21} k^3 = 2 \cdot \left(\frac{21 \cdot 22}{2}\right)^2 = 2 \cdot 231^2 = 2 \cdot 53361 = 106722$

\textit{Erklärung: $\frac{21 \cdot 22}{2} = 231$, dann $231^2 = 53361$.}

\bigskip

\textbf{Summe 2:} $\displaystyle 2 \sum_{k=1}^{21} k = 2 \cdot \frac{21 \cdot 22}{2} = 2 \cdot 231 = 462$

\bigskip

\textbf{Summe 3:} $\displaystyle 17 \sum_{k=1}^{21} 1 = 17 \cdot 21 = 357$

\bigskip

\textbf{Ergebnis:}
\[
    106722 - 462 + 357 = 106617
\]

\[
    \boxed{\sum_{k=1}^{21} \left(2k^3 - 2k + 17\right) = 106617}
\]

\end{document}
